% Part: lambda-calculus
% Chapter: lambda-definability
% Section: lists

\documentclass[../../../include/open-logic-section]{subfiles}

\begin{document}

\olfileid{lam}{ldf}{lst}
\olsection{فهرست‌ها}

کدگذاری فهرست $[M_1, M_2, \ldots, M_n]$ را چنین تعریف می‌کنیم:
\[
  [M_1, M_2, \ldots, M_n] \ident \lambd[sz][s M_1 (s M_2 (\ldots (s M_n z)))]
\]
به‌طور غیررسمی، ترمی که فهرست را نمایش می‌دهد تابع «انباشتگر» آن
فهرست است: تابعی که دو ترم می‌پذیرد. ترم نخست تابعی است که مقدار
انباشته‌شده را همراه با یک عنصر تازه می‌گیرد و مقدار انباشته‌شدهٔ تازه
را بازمی‌گرداند؛ ترم دوم مقدار انباشته‌شدهٔ آغازین است. این تابع عناصر
را یکی‌یکی می‌انبازد و سرانجام نتیجه را بازمی‌گرداند.

\begin{digress}
  اگر با برنامه‌نویسی تابعی آشنا باشید، این تعریف را شبیه
  \emph{تاخوردگی راست} خواهید یافت: فهرست به‌صورت تابع تاخوردگی راست
  عناصر درون آن تعریف می‌شود.
\end{digress}

بر پایهٔ این کدگذاری می‌توانیم چند تابع سودمند را به‌آسانی تعریف کنیم:
\begin{align*}
  \fn{Sum} & \ident
  \lambd[l][l \, (\lambd[xa][\fn{Add} \, x \, a])\,  \num{0}]\\
  \fn{Len} & \ident
  \lambd[l][l \, (\lambd[xa][\fn{Add} \, a \, \num{1}]) \num{0}]
\end{align*}

$\fn{Sum}$ مجموع فهرستی از اعداد چرچ را محاسبه می‌کند. این تابع با
انجام انباشت روی فهرست کار می‌کند: مقدار آغازین $\num{0}$ است و برای
هر عنصر، $\fn{Add} \, x\, a$ را به‌عنوان مقدار تازه بازمی‌گرداند، که
در آن $x$ عنصر کنونی و $a$ مقدار انباشته‌شده است. نتیجه، مجموع همهٔ
عناصر است.

\begin{prob}
  $\fn{Len}$ چه می‌کند؟ توضیح دهید.
\end{prob}

\begin{prob}
  تابعی تعریف کنید که فهرست را معکوس کند.
\end{prob}
\end{document}

